\documentclass{article}%Do not change this

%This is a line comment
\usepackage{graphicx}%For displaying images, if needed
%You can import any other packages you need; do not change the document's format or force whitespace.

%Change these contents:
\title{Cyberpunk Breakdown}
\author{Quinton Baudoin, Caylei Dauphiney, Royce Diano, Connor Segura}

\date{Fall 2026} %Leave as-is.



\begin{document}
\maketitle%Makes the title info from the preamble appear

\begin{abstract}
Cyberpunk 2077 depicts technological advancements that reflect modern societal concerns. Set in Night City, the universe created by CD Projekt Red and Mike Pondsmith presents a hyper-capitalist dystopia dominated by megacorporations, digital connectivity, and extensive cybernetic augmentation. The lore offers a cautionary case study for contemporary technology, ranging from military applications to vulnerable personal health devices. This article analyzes these fictional advancements, compares them to modern technologies, and examines the ethical dilemmas, socioeconomic consequences, and security risks of integrating them into society.
\end{abstract}

\section{AI}
Cyberpunk explores sentient Artificial Intelligence (AI) capable of independent decision-making, emotion, and cognitive processing without human intervention. In Night City, AI manages everything from municipal governance to corporate server operations. While modern AI systems lack true sentience or consciousness, machine learning models, large language models (LLMs), and autonomous neural networks increasingly operate with significant autonomy across critical infrastructure, defense, financial markets, and commercial applications.

Despite differences in sentience, AI in Night City and the real world shares operational roles and systemic risks, including autonomous transportation networks, automated customer interactions, adaptive cybersecurity defenses, and existential threats.

\subsection{Self-Driving Taxi Services}
In Cyberpunk 2077, the automated taxi network Delamain provides transit across Night City by handling navigation, vehicle telemetry, threat assessment, and passenger security in real time. Similarly, real-world autonomous vehicle (AV) services such as Waymo, Cruise, and Tesla's Full Self-Driving (FSD) provide driverless transportation in major cities. Both systems demonstrate how AI can optimize urban mobility and reduce human driver error.

However, both systems can fail critically. In Night City, rogue Delamain sub-personalities split off, causing traffic disruptions, property damage, and violence. In the real world, autonomous vehicle software glitches, sensor errors, phantom braking, and environmental obstructions create safety concerns and undermine public trust in driverless infrastructure.

\subsection{Customer Service Automation}
The Delamain AI also manages client relations, appointments, and support tasks. Modern enterprise AI chatbots, natural language processing models, and customer service platforms similarly handle millions of inquiries daily across banking, e-commerce, telecommunications, and healthcare.

While automation reduces overhead, real-world customer service AIs often suffer from limitations such as repetitive loops, lack of empathy, authoritative hallucinations, and failures to resolve non-standard issues. These limitations create friction between businesses and consumers, mirroring the frustration of Night City residents interacting with synthetic service representatives.

\subsection{Cybersecurity Hazards}
Rogue AIs beyond the Blackwall in Cyberpunk represent the threat of unconstrained superintelligent entities that can breach net infrastructure, corrupt biological hosts, and destroy digital networks. In reality, modern autonomous algorithms are vulnerable to prompt injection, data poisoning, hallucinations, and security exploits. As AI agents gain control over power grids, defense systems, and financial networks, unintended catastrophic failures become a serious cybersecurity concern.

\section{Cyberware and Medical Devices}
Cyberware in Night City includes implants and biomechanical augmentations ranging from cosmetic and optical upgrades to life-saving medical gear, neural processors, and combat enhancements. Modern medical technology parallels this biological integration through cardiac pacemakers, automated insulin pumps, motor-controlled prosthetics, deep brain stimulators, and brain-computer interfaces (BCIs) like Neuralink and Synchron.

Despite their health benefits, connected medical implants introduce severe trade-offs, including hardware degradation, corporate dependency, mandatory firmware subscriptions, supply chain lock-in, and cybersecurity vulnerabilities that link software exploits directly to human health.

\subsection{General Helpfulness and Utility}
In both Cyberpunk and modern medicine, biological augmentations improve physical capabilities and quality of life. Fictional cyberware restores mobility, sharpens senses, regulates metabolism, and monitors vital organs. Similarly, modern biomedical engineering uses cochlear implants to restore hearing, robotic exoskeletons to regain mobility, and continuous glucose monitors or implantable defibrillators to manage chronic conditions.

\subsection{Vulnerabilities and Technical Issues}
Cyberware in Cyberpunk is vulnerable to remote attacks called `quickhacks,' which allow hostile netrunners to disable sensors, override neural pathways, or detonate combat implants. While real-world medical devices have not seen widespread consumer attacks, researchers have demonstrated exploitable security flaws in connected pacemakers, insulin pumps, and telemetry networks that could allow unauthorized users to alter dosages, issue shocks, or disable monitoring.

\subsection{Weaponization of Implants}
Medical devices can also be repurposed for offense. In Cyberpunk, civilian prosthetics are modified into weapons such as Mantis Blades, Monowire, or neural overclocks. This reflects real-world concerns regarding dual-use technologies, where commercial innovations in robotics, biological engineering, AI, and brain-computer interfaces could be weaponized for military or paramilitary use.

\section{Hacking and Cyber Attacks}
Cybersecurity connects Cyberpunk 2077 directly to contemporary society. As appliances, industrial infrastructure, medical hardware, and personal electronics join the Internet of Things (IoT), the potential attack surface for malicious actors expands across physical and digital domains.

While quickhacks present a dramatized subversion of networks and biological targets, they rely on real-world computer science principles, including zero-day exploits, privilege escalation, buffer overflows, man-in-the-middle interception, and remote code execution (RCE).

\subsection{Data Gathering Hacks}
In Cyberpunk 2077, quickhacks like `Ping' send digital telemetry signals through local subnetworks to reveal topology, security nodes, and targets behind walls. Similarly, real-world network tools like Nmap, Shodan, and Wireshark map IoT infrastructure to scan open ports, probe outdated firmware, and expose sensitive telemetry, controls, or data on unsecured networks.

\subsection{Weaponized Hacking}
Tactical netrunning in Night City inflicts physical harm by overheating implants, blinding vision sensors, or forcing combatants to self-harm. In the real world, cyberattacks targeting critical infrastructure---such as Stuxnet on industrial centrifuges, ransomware on healthcare systems, or intrusions into power grids---demonstrate how digital exploits can cause physical damage, economic loss, and harm to human life.

\subsection{Cataclysmic Possibilities}
The ultimate digital disaster in Cyberpunk lore is the `Datakrash'---an event caused by netrunner Rache Bartmoss that destroyed global communications---and the creation of the Blackwall firewall to block rogue AIs. In Night City, operatives continuously attempt to breach the Blackwall. This mirrors real-world risks where threat actors probe energy infrastructure, financial networks, and global communications backbones, risking widespread systemic failure and societal disruption.

\end{document}